\documentclass[11pt]{article}
\usepackage[T1]{fontenc}
\usepackage{amsmath,amssymb,amsthm}
\usepackage[margin=1in]{geometry}
\usepackage[hidelinks]{hyperref}
\newtheorem{theorem}{Theorem}
\newtheorem{lemma}[theorem]{Lemma}
\newtheorem{proposition}[theorem]{Proposition}
\newtheorem{corollary}[theorem]{Corollary}
\newcommand{\C}{\mathcal C}
\newcommand{\D}{\mathcal D}
\newcommand{\aa}{\mathfrak a}
\newcommand{\RR}{\mathbb R}
\title{Geometric fibre maxima over normal targets\\
Finite-cover envelope descent and canonical norm bounds}
\author{Independent scoped ORIGINAL P12.3 deductions; unreviewed draft}
\date{9 October 2026}
\begin{document}
\maketitle
\begin{abstract}
For a finite domain extension over an excellent normal local domain, we
prove that the maximum over the entire restriction fibre preserves the
independently geometric continuous and decreasing-completed cones.
Characteristic-polynomial coefficient ideals and actual primary
truncations give the maximum of a primary profile as a finite maximum
of actual primary profiles, including every nonminimal support and
inseparable multiplicity. Consequently raw envelopes of pulled target
obstacles pull back exactly and are fibre-constant even before regularity
is known. Greatest minorants exist downstairs exactly when the invariant
source obstacle has one. Any actually known source continuity theorem
therefore descends along a finite cover of a normal target, without a
group action or flatness. Generic-degree norm traces give every-degree
canonical error bounds and a model rate. This does not assert that every
normal singular germ admits a finite regular cover.
\end{abstract}

\section{Exact domain and the full source-qualified trace}
Let $R$ be an excellent normal local domain of positive dimension, and
let $R\subset S$ be a finite domain extension. Regard $S$ as semilocal
with all closed centres above the maximal ideal of $R$. Fix a proper
primary $\aa\subset R$, and put
\[
 X=V_\aa(R),\qquad Y=V_{\aa S}(S),\qquad \rho:Y\to X,\qquad
 K=\operatorname{Frac}(R),\quad L=\operatorname{Frac}(S),\quad n=[L:K].
\]
All prime supports, infinite individual values, and the full compact
pointwise topology are retained. The source classes are the finite product
classes over its marked centres. They agree with semilocal ideal profiles
whose ideals contain a power of $\aa S$.

The sole incoming correctness grant is ENTIRE qualified G331, ALL FIVE
fields and every model, support, source, radial and historical-audit
qualification. Positivity means actual negative effective finite real
Cartier presentations on directed ACTUAL primary blowups, with degree
at least zero on EVERY contracted proper integral curve before comparison
or limits. G331's independently proved comparison supplies rational
primary-profile density, finite maxima and the actual decreasing-net
completion $\D$ of its continuous cone $\C$.

We use the complete normal-target coefficient trace of G331 n3 at its
precise domain. Choose a finite field $M/K$ containing all embedding
images of $L$, and the finite integral closure $E$ of $R$ in $M$.
For every $v\in X$, lift it to $E$ by quotient-support lifting.
Compose with the embeddings to obtain source points $w_\sigma$.
If $q$ is the generic inseparable length, the measure
\[
 \mu_v=q\sum_{\sigma:L\hookrightarrow M}\delta_{w_\sigma}
\]
is independent of that lift, has mass $n$, and has support the ENTIRE
restriction fibre. Thus every fibre is nonempty and finite, with at most
$n$ points. Each atom has positive integer weight. The trace
$T(H)(v)=\int H\,d\mu_v$ is continuous for every continuous $H$,
has norm $n$, and sends a negative rational source profile to a negative
norm-coefficient profile downstairs. This theorem uses Gauss values in
the fraction field of the support quotient, compact gauge-lattice density,
and integral normal-target coefficients. It assumes no defectless
ramification formula and no extension of a nonzero-support semivaluation
to $L$. Its full qualifications remain in force.

\section{Characteristic coefficients on every support}
Let $J\subset S$ contain $(\aa S)^N$ for an integer $N\ge1$.
Choose finite $R$-module generators $b_1,\ldots,b_s$ of $J$ and set
$b(T)=\sum_jT_jb_j$. Its generic characteristic polynomial is
\[
 P_J(Z,T)=Z^n+\sum_{\ell=1}^n(-1)^\ell c_\ell(T)Z^{n-\ell}
         =\prod_{\sigma}(Z-\sigma b(T))^q.
\]
Let $C_\ell(J)$ be the ideal generated by the coefficients of
$c_\ell(T)$. All these coefficients belong to $R$: each is integral over
$R$ as a polynomial expression in integral conjugates, and belongs to
$K$ by the characteristic-polynomial construction. Normality places it
in $R$. Changes of $R$-module generators make linear substitutions over
$R$; substitutions both ways prove generator independence.
Some $C_\ell$ can be zero or nonprimary, particularly in positive
characteristic. Define the actual unit-or-primary ideals
\[
 D_\ell(J)=C_\ell(J)+\aa^{N\ell}.
\]
Adding this boundary power is essential.

\begin{lemma}\label{coeff}
For every target semivaluation, with arbitrary support,
\[
 \min_{\rho(w)=v}w(J)
   =\min_{1\le\ell\le n}\frac{v(C_\ell(J))}{\ell}
   =\min_{1\le\ell\le n}\frac{v(D_\ell(J))}{\ell}.
\]
Zero coefficient ideals have value infinity before truncation.
\end{lemma}
\begin{proof}
Lift $v$ to $E$ as above, kill its support first, and enter the resulting
valued quotient fraction field. On its rational-function extension in
$T$, use the Gauss valuation, the minimum coefficient value. This is
multiplicative because normalizing least-valued coefficients and reducing
gives nonzero residue polynomials with nonzero product.
Each root $\sigma b(T)$ of the specialized polynomial has finite value
\[
 \alpha_\sigma=\min_j w_\sigma(b_j)=w_\sigma(J)\le N.
\]
The equality uses the finite $R$-module generating set; coefficients in
$R$ have nonnegative values. Finiteness uses the boundary power in $J$.
Each root is counted $q$ times.

Put $a=\min_\sigma\alpha_\sigma$ and let $r$ be the number of roots,
with multiplicity, having value exactly $a$. Every elementary symmetric
coefficient of order $\ell$ has value at least $\ell a$.
For order $r$, the product of ALL $r$ least-valued roots is the unique
summand of value $ra$. Every other $r$-element subset includes a root of
strictly larger value, so its value is larger. The unique lowest summand
cannot cancel. Hence the order-$r$ coefficient has value exactly $ra$.
This works with repeated roots and in every characteristic; no binomial
coefficient nonvanishing premise is used.
The Gauss value of $c_\ell(T)$ is precisely $v(C_\ell(J))$.
G331's full-support measure theorem identifies the roots' values with
values on the ENTIRE fibre, proving the first equality.
Finally $a\le N$ and
$v(D_\ell)/\ell=\min(v(C_\ell)/\ell,N)$, so truncation does not change
the minimum. All root evaluations took place after killing the support.
\end{proof}

\section{Positive geometric fibre maxima}
Define the full-fibre maximum by
\[
 M_\rho u(v)=\max_{\rho(w)=v}u(w).
\]
For finite continuous $u$ it is continuous. One proof using G331's trace
makes the topology explicit: choose $c$ so $B=u+c\ge0$. For integer $p\ge1$,
\[
 M_\rho B\le(T(B^p))^{1/p}\le n^{1/p}M_\rho B.
\]
The lower bound uses positive integer atom weights; the upper uses mass
$n$. The middle expression is continuous, and the uniform error is at
most $(n^{1/p}-1)\|B\|_\infty\to0$. Subtract $c$.
This ordinary continuity argument does not establish geometric positivity.

\begin{theorem}\label{max}
The operator $M_\rho$ sends $\C_S$ to $\C_R$ and $\D_S$ to $\D_R$.
It is monotone, is a contraction on finite functions, and is left adjoint
to pullback in the pointwise order:
\[
 M_\rho u\le h\quad\Longleftrightarrow\quad u\le\rho^*h.
\]
The source operator $\rho^*M_\rho$ is the least fibre-constant majorant,
and is idempotent on these geometric classes.
\end{theorem}
\begin{proof}
For a rational source primary/unit profile, Lemma \ref{coeff} gives
\[
 M_\rho(-w(J)/k)
   =\max_{1\le\ell\le n}
          \left(-\frac{v(D_\ell(J))}{\ell k}\right).
\]
Every proper $D_\ell$ has an ACTUAL primary blowup and negative effective
relatively ample presentation, with positive degree on every contracted
integral curve before any limit. Units give zero. The finite maximum
belongs to $\C_R$ by G331's independently geometric comparison and maximum
stability. Rational-profile density and the elementary fibre-maximum
supremum-norm contraction extend this to every $\C_S$ function.

If $u=\inf_i u_i$ is an actual decreasing $\C_S$ net, the fibre is finite.
Consequently
\[
 \max_{\rho(w)=v}\inf_i u_i(w)=\inf_i\max_{\rho(w)=v}u_i(w).
\]
For a finite limit, choose a common late index approximating all finitely
many pointwise limits from above; bottom coordinates can be made below
that same finite threshold. If all limits are bottom, a common late index
makes every value below any prescribed negative number. This proves the
identity including bottom. The right side is an ACTUAL decreasing
$\C_R$ net, so belongs to $\D_R$.
Adjunction, the least-majorant statement and idempotence follow pointwise
from nonempty full fibres. This is a positive operator theorem, not a
designation of arbitrary fibre extrema as psh.
\end{proof}

\section{Exact envelope transfer without a group}
For a finite continuous signed $f$ on $X$ write $p_R(f)$ for the raw
supremum over all geometric minorants, separately in $\C$ and $\D$.
The canonical argument below makes these two suprema equal and finite.
\begin{theorem}\label{env}
On EVERY source point,
\[
 \rho^*p_R(f)=p_S(\rho^*f).
\]
The formula holds for each raw class independently. The source raw
envelope is therefore full-fibre-constant, even before continuity is known.
A greatest $\C$ or $\D$ minorant exists on the target if and only if the
corresponding pulled source obstacle has one. Universal source continuous
envelopes imply universal target continuous envelopes.
\end{theorem}
\begin{proof}
Pullback preserves the geometric classes and gives one raw inequality.
Any source competitor $u\le\rho^*f$ has, by Theorem \ref{max},
a target competitor $M_\rho u\le f$ and
$u\le\rho^*M_\rho u\le\rho^*p_R(f)$. Take its raw supremum.
This proves equality without assuming flatness, a group action, or
positivity of an increasing supremum.

If the target has a greatest minorant, its pullback equals the source raw
supremum and is the source greatest minorant.
If the source has a greatest minorant, the proved formula makes it
fibre-constant; its maximum descends in the same positive class and is
the target greatest minorant. The canonical argument gives continuity
and equality of the greatest $\C,\D$ members whenever either exists.
The converse universal statement concerns only source pulled obstacles,
not arbitrary source obstacles.
\end{proof}

For any finite continuous signed source obstacle $F$, ordinary $M_\rho F$
is continuous by the trace argument, and one also has the raw Jensen bound
\[
 M_\rho p_S(F)\le p_R(M_\rho F).
\]
Indeed $M_\rho$ carries every source competitor to a target competitor;
suprema over competitors commute with the finite fibre maximum.
No equality or strictness criterion for arbitrary $F$ is asserted.
In a tower of finite domain extensions with normal targets and intermediate
ring, full fibre maxima compose exactly: a total fibre is the union of the
last-stage fibres over the intermediate fibre. The same elementary identity
applies to $\C,\D$ by their proved preservation, using extended base gauges.

\section{Canonical norm bounds at every degree}
For either ring define $A_m(f)=\{b:v(b)+mf(v)\ge0\ \forall v\}$ and
$h_m=-v(A_m)/m$. Ultrametricity and a sufficiently high boundary power
make $A_m$ unit-or-primary; $A_mA_l\subset A_{m+l}$.
Each $h_m$ is an actual ample-primary geometric minorant.
Rational approximation of any $\C$ competitor $g$ by $-v(I)/k$ and a
negative $\epsilon$ gauge correction gives
$I^N\aa^{\lceil Nk\epsilon\rceil}\subset A_{Nk}$.
Thus the raw $\C$ supremum is $\sup_m h_m$.
For an actual $\D$ minorant $\inf_i g_i$, the open strict tests
$\{g_i<f+\epsilon\}$ cover the compact section; one directed index works
everywhere, making $g_i-\epsilon$ a $\C$ minorant.
This proves equality of both raw suprema, including partial bottom.

For $N=km+r$, $0\le r<m$, grading gives
$h_N\ge(km/N)h_m-r\|v(A_1)\|_\infty/N$.
Every degree therefore tends pointwise to the raw supremum.
If it is continuous, a finite strict-test cover and a common multiple
of its indices give uniform convergence by this same inequality.
The uniform limit lies in $\C$ and is the greatest minorant.
A greatest $\D$ minorant is both the lower semicontinuous raw supremum
and an upper semicontinuous decreasing completion, hence continuous.
Negative gauge shifts give the $1$-Lipschitz obstacle bound.
These arguments require no finite generation.

Put $J_m=A_m^S(\rho^*f)$, $I_m=A_m^R(f)$. Full surjectivity gives
$I_m=J_m\cap R$. Let $\mathcal N(J_m)$ be G331's generic-degree
norm-coefficient ideal. Its all-support trace identity gives
\[
 v(\mathcal N(J_m))
   =\int w(J_m)\,d\mu_v\ge-nmf(v),
 \qquad \mathcal N(J_m)\subset I_{nm}.
\]
The inclusion follows by the canonical tests on EVERY coefficient and
EVERY target point; no false assertion that conjugates belong to $J_m$
is needed. If $p_S(\rho^*f)$ is continuous, let
$\beta_m=\|p_S-h_m^S\|_\infty$. It tends to zero by the canonical argument.
Since $p_S=\rho^*p_R$, normalized trace yields
\[
 h_{nm}^R\ge \frac1nT(h_m^S)\ge p_R-\beta_m.
\]
For every $N\ge n$, write $N=nm+r$, $0\le r<n$. Grading and $p_R\le0$ give
\[
 0\le p_R-h_N^R
 \le\frac{nm}{N}\beta_m+\frac rN\|v(I_1)\|_\infty
 \le\beta_{\lfloor N/n\rfloor}
       +\frac{n-1}{N}\|v(I_1)\|_\infty .
\]
If the exact source model estimate is $\beta_m\le K_0/m$, the sharper
first bound gives $(nK_0+(n-1)\|v(I_1)\|_\infty)/N$ for $N\ge n$.
For $N<n$, $I_1^N\subset I_N$ bounds the error by $\|v(I_1)\|_\infty$.
Thus $n(K_0+\|v(I_1)\|_\infty)/N$ is a single all-degree model bound.
No source model estimate is inferred for arbitrary continuous data.

\section{Proved sectors, sources and remaining original scope}
Whenever all marked source centres are smooth affine CLOSED germs over
an algebraically closed characteristic-zero field, their finite product
has universal source continuity by the exact qualified BFJ theorem.
Theorem \ref{env} proves target continuity for every continuous signed
obstacle, with no finite-group presentation required. The same statement
applies to actually supplied regular formal power-series source germs
through G331's completion homeomorphism and exact primary contraction.
In any characteristic or mixed characteristic the abstract transfer theorem
remains valid; it supplies continuity only when the particular source
continuity theorem has actually been established.

ENTIRE G331 n3 was freshly physically reread for its generic inseparable
lengths, normal-target integrality, all-support quotient Gauss proof,
full-fibre support and continuous trace. The current430 canonical and smooth
proof and the six complete immediate BFJ units retain the fresh current
ROOT430 reading and rejudged applications, with no extra primary reading
credit. The smooth domain, real $t=m+C$, thinness bounded on a FIXED
complex, substantive closed-simplex continuity and ALL retractions remain
mandatory. Historical principalization, multiplier subadditivity, Izumi
and conventional finite/excellence algebra are disclosed background;
the deeper recursive source/book/pixel audit is incomplete.
The characteristic coefficients, minimal-root argument, exact primary
truncation, positive maxima, full envelopes and norm bounds are proved here.
No new external valuation-transitivity theorem is invoked.

Other primary gauges divide every homogeneous function by the same
bounded positive continuous denominator, constant on each full fibre.
Maxima, envelopes, trace normalization and canonical tests transport
with that one factor; canonical ideals stay unchanged for the transported
obstacle. G331 completion carries proved target continuity by its actual
boundary/cone homeomorphism; no arbitrary model-category equivalence or
new formal-invariant identity is asserted. Empty sections are direct.

Normality of the target is essential to the coefficient construction.
No version for arbitrary nonnormal conductor targets is claimed.
Source universal continuity is not inferred from target continuity.
No arbitrary normal singular germ is assumed to have a finite regular
cover. General excellent NORMAL dimension at least three and unrestricted
positive/mixed-characteristic higher-dimensional regular-source continuity,
intrinsic conductor classification, and full original P12.3 remain OPEN.
Products, energy, full operators and literal radial comparison stay separate.
This proof was developed after the native298 launch without reading its
active mathematics or feeding it back, and before any native297 mathematical
inspection. Compilation and acceptance are unverified; candidate is false.
\end{document}
